% P8 experimental setup: honest description of the custom dataset + both arms
\section{Experimental Setup}
\label{sec:p8-setup}

\paragraph{Dataset (custom, synthetic, single-configuration).}
We are explicit up front: the benchmark is a \emph{custom synthetic dataset},
not a public benchmark and not production data from any institution. It is
generated with scikit-learn's \texttt{make\_classification}
\citep{pedregosa2011sklearn} with a fixed seed (42): $50{,}000$ rows, $20$
anonymized numeric features $V_1,\dots,V_{20}$ (10 informative, 2 redundant),
two clusters per class, class separation $0.8$, $1\%$ label noise
(\texttt{flip\_y}$=0.01$), and a $1\%$ positive (``fraud'') rate---the
class-imbalance regime typical of card fraud \citep{dalpozzolo2018realistic}.
The anonymized $V$-feature convention deliberately mirrors the PCA-transformed
feature sets that single-institution card-fraud datasets expose after
confidentiality processing; the reader should treat our data as a stylized
stand-in for one institution's post-processing feature view, with none of the
temporal drift, verification latency, or covariate shift of real transaction
streams \citep{dalpozzolo2018realistic}. Four aggregate features are appended
per row (mean, standard deviation, max, min of $V_1,\dots,V_{20}$), giving 24
features total.

\paragraph{Split and metric.}
An 80/20 stratified train/test split (seed 42) yields $40{,}000$ training and
$10{,}000$ test rows, with roughly $100$ positives in the test set. The primary
metric is held-out ROC-AUC; F1, precision, and recall at the default threshold
are logged alongside. Both arms consume \emph{identical} splits: the split
files are written to disk once and the LLM arm is fine-tuned on a textual
serialization of the same training rows.

\paragraph{Scorer arm: XGBoost.}
The tree baseline is a stock XGBoost classifier \citep{chen2016xgboost} with
$200$ estimators, maximum depth $6$, learning rate $0.05$, subsample and
column-subsample $0.8$, and \texttt{scale\_pos\_weight}$=7$ to counter the
$99{:}1$ imbalance; the evaluation objective is AUC. No hyperparameter search
was run beyond this single sensible configuration---the point of the baseline
is what an ordinary, lightly tuned tree achieves, not a leaderboard entry. The
full implementation is a single short script (\texttt{train\_xgboost.py},
released with the paper) that also emits wall-clock training and per-10k-row
inference times. One disclosure belongs here rather than in the limitations:
the current quick artifact reports test AUC $0.7955$ for this XGBoost arm
on $10{,}000$ held-out rows
(\texttt{experiments/results/quick\_20260704/qp8\_fraud.tsv}). An older
21~June~2026 internal record listed XGBoost at AUC $0.975$, but we could not
reconstruct that environment and do not use that number as a reproducible
headline result (\secref{sec:p8-limitations}).

\paragraph{Challenger arm: fine-tuned LLM.}
The current challenger is a Qwen3.5-4B SFT run fine-tuned on row-to-text
serializations of the training split, in the style of TabLLM
\citep{hegselmann2023tabllm}: each row is rendered as a feature-name/value
string and the model is trained to emit a fraud/legitimate judgment. Because
the natural-rate test set contains too few positives for a stable LLM AUC in
the quick run, the evaluation uses a 500-row positive-enriched held-out subset
with 20\% fraud, disjoint from training
(\texttt{qp8-fraud-sft\_manifest.json}). The final quick artifact reports
accuracy $0.792$ and AUC $0.48268$ (rounded to $0.4827$ in
\texttt{qp8-fraud-sft.tsv}; summarized in \texttt{qp8\_fraud.tsv}). An older
21~June~2026 internal record listed a fine-tuned LLM at AUC $0.948$, but no
released artifact reproduces it, so we retain it only as historical provenance
(\secref{sec:p8-limitations}).

\paragraph{What this setup can and cannot support.}
The synthetic generator gives us a clean, reproducible, shareable testbed with
a known imbalance and noise level, and it suffices for the paper's actual
claim: a comparison of \emph{roles} (who should hold the scorer seat, and what
the LLM is uniquely positioned to do elsewhere in the pipeline). It cannot
support absolute performance claims about production fraud systems, and no
number in this paper should be quoted as an expected production AUC. Where we
report operational quantities that depend on facts outside the dataset---
analyst costs, latency envelopes---we scope them explicitly as internal
estimates.
